% Generated by tools/edn_latex.py from reports/edn.md.
% Do not edit: change reports/edn.md and run `make edn`.
\ednentry{
  id = {EDN-42},
  title = {A categorical value the training rows never saw becomes missing, and the holdout's country is allowed to vanish},
  date = {2026-09-30},
  milestone = {M2: Reproducibility (feature encoding; the holdout feeds M6: Monitoring)},
  topic = {Feature Encoding, Monitoring},
  participants = {@lukas2510},
  decision = {Every categorical feature leaves the \texttt{features} stage as a \texttt{category} over the levels the training split holds, for every split and for a request alike. A value outside those levels becomes missing rather than a level of its own. The accepted consequence is that the \texttt{ES} holdout's \texttt{country\_\allowbreak{}code} is missing in \textbf{every row} of its feature matrix, because \texttt{ES} is held out by construction and is therefore not a training level. This is documented rather than worked around: the stage warns, naming any column a written split never fills, a test pins the holdout case as intended, and the problem specification states it.},
  rationale = {\emph{Alternatives considered:}
\begin{itemize}[leftmargin=*, topsep=2pt]
\item \textbf{Option A (chosen): training levels only, an unseen value is missing, and the holdout's empty column is documented.}
\newline Pros: a code means the same car in training, in test, in the holdout and at request time, which is the only way a fitted model can be applied to anything but its training frame. It is exactly what \hyperref[EDN-18]{EDN-18} already decided for the serving path -- an unseen country is accepted and passed to the model as unknown -- so the pipeline and the API agree without a second rule. The \texttt{ES} matrix being country-blind is the new-market scenario stated honestly: the model has never seen that market, which is the whole point of \hyperref[EDN-03]{EDN-03}. Costs nothing to implement.
\newline Cons: a column that is empty in every row looks exactly like a defect, and the holdout is not replayed until M6, months after anyone remembers why. Mitigated by the warning, the test and the specification sentence rather than by changing the encoding.
\item \textbf{Option B: give the holdout country a level of its own.}
\newline Pros: the column is no longer empty, so nothing looks broken.
\newline Cons: hands the model a code it was never fitted on, which is worse than missing: a tree splits on it with no training example behind the branch. It also renumbers every other level relative to the fitted model unless the level is appended, and appending a level nothing was fitted on is a contract the artefact cannot honour. And it would make the drift scenario measure something other than a new market.
\item \textbf{Option C: drop \texttt{country\_\allowbreak{}code} from the feature sets, so the holdout matrix has no such column.}
\newline Pros: no empty column anywhere.
\newline Cons: throws away a feature the problem specification names for every other split, where it is filled in 99.98\,\% of rows, in order to tidy up one holdout frame. Country is a real price driver across seven markets.
\item \textbf{Option D: let each frame decide its own levels.}
\newline Pros: no frame ever has an empty categorical.
\newline Cons: the failure this whole mechanism exists to prevent. Every split would get different codes, so a model fitted on \texttt{train} could not be applied to \texttt{test} at all, and a request could not be scored. It is also a leak: validation, test and the holdout would be deciding the feature space.
\end{itemize}
\par
\emph{Why:} The encoding rule was effectively settled by EDN-02 (native categoricals) and EDN-18 (accept an unseen value, do not reject it); what was open is whether the holdout's fully empty \texttt{country\_\allowbreak{}code} is a defect to fix or a consequence to document. It is a consequence, and the two ways of ``fixing'' it are both worse: option B feeds the model an unfitted code and option C deletes a real feature. The actual defect was that nothing said so -- the reasoning existed only in a pull-request description, which does not survive a squash merge -- so the fix is the warning, the test and one sentence in the specification. Whether the model should instead be fitted and gated only on the makes and markets the API serves is a separate question that is deliberately not settled here.},
  ai = {Alternative assessment, Recommendation},
  response = {Accepted},
  assessment = {An adversarial review found the empty column by running the stage on the real snapshot and checking every column of every split, and established that \texttt{country\_\allowbreak{}code} is the only column fully null in the holdout and not also fully null in training. It read the situation as faithful to the drift scenario rather than as a bug, and recommended documenting it in three places rather than changing the encoding, explicitly declining to drop the column or to add a level because those are modelling changes nobody had asked for. That framing is what made this an entry rather than a patch.},
  aievidence = {Claude Code session on 2026-09-30, adversarial review of \href{https://github.com/mlops-2627q1-mds-upc/MLOPS_Recommenditos/pull/55}{PR \#55}: ``this is faithful to the drift scenario rather than a defect, because \texttt{ES} is held out by construction, so the model genuinely has never seen that market, and EDN-18 already requires an unseen value to be accepted rather than rejected. What is unacceptable is that nothing says so.''},
  evidence = {\href{https://github.com/mlops-2627q1-mds-upc/MLOPS_Recommenditos/blob/main/recommenditos/data/build_features.py}{\texttt{recommenditos/\allowbreak{}data/\allowbreak{}build\_\allowbreak{}features.\allowbreak{}py}} (\texttt{\_\allowbreak{}warn\_\allowbreak{}about\_\allowbreak{}unobserved\_\allowbreak{}columns}); \href{https://github.com/mlops-2627q1-mds-upc/MLOPS_Recommenditos/blob/main/tests/test_features.py}{\texttt{tests/\allowbreak{}test\_\allowbreak{}features.\allowbreak{}py}} (\texttt{test\_\allowbreak{}the\_\allowbreak{}holdout\_\allowbreak{}country\_\allowbreak{}is\_\allowbreak{}missing\_\allowbreak{}throughout\_\allowbreak{}its\_\allowbreak{}matrix\_\allowbreak{}on\_\allowbreak{}purpose}); \href{https://github.com/mlops-2627q1-mds-upc/MLOPS_Recommenditos/blob/main/docs/docs/problem-spec.md}{problem specification} section 2; \hyperref[EDN-03]{EDN-03}; \hyperref[EDN-18]{EDN-18}; \href{https://github.com/mlops-2627q1-mds-upc/MLOPS_Recommenditos/pull/55}{PR \#55}.},
}
